\chapter{Diskussion}
\label{chap:diskussion}

\section{Interpretation der Ergebnisse}

\subsection*{F1 — Zeitreduktion}

Die starke Variation der Zeitreduktion (9,4~\%–99,5~\%) liegt nicht am Genre --- der Kruskal-Wallis-Test zeigt keinen signifikanten Genreeffekt auf die jährlichen Verbesserungsraten. Entscheidender sind spielspezifische Faktoren: Spiele mit kurzer Grundlaufzeit und nutzbaren Programmfehlern (Minecraft~< 2~Min., Portal~< 10~Min.) lassen sich stärker optimieren als Spiele, bei denen physikalische oder spielbedingte Grenzen eine weitere Verkürzung verhindern (Super Mario Bros.: jedes Frame zählt, wenige Glitches verfügbar). Innerhalb des RPG-Genres ist die Spiellänge der begrenzende Faktor --- Final Fantasy~VII dauert im Normalspiel über 35~Stunden.

\subsection*{F2 — Sättigung}

Der universelle Strukturbruch (alle 17~Spiele, Chow-Test signifikant) ist das stärkste Ergebnis von F2. Er deutet darauf hin, dass Speedrunning-Optimierung nicht linear verläuft: Eine frühe \emph{Hochaktivitätsphase} (intensive Routing-Suche, Entdeckung spielverändernder Glitches) geht in eine flachere \emph{Plateau-Phase} über, in der nur noch marginale Verbesserungen möglich sind. Die Tatsache, dass der Strukturbruch häufig schon bei den ersten 10~Weltrekorden auftritt, legt nahe, dass frühe Mitglieder der Community den größten Teil der erreichbaren Optimierung erzielen. Das hohe $R^2$ aller Modelle ($> 0{,}70$) bestätigt, dass der Kurvenverlauf statistisch gut beschreibbar ist, unabhängig vom gewählten Modelltyp.

\subsection*{F3a — Post-Breakthrough-Dynamik}

Der hier gewählte Ansatz unterscheidet sich vom Chow-Test-Bruchpunkt aus F2: Der Chow-Test findet den Punkt mit dem stärksten Steigungswechsel in der \emph{kumulierten} Zeitreihe und liegt typischerweise bei den ersten 10~WRs, weil die Verbesserungsrate dort am steilsten abfällt. Der Breakthrough hingegen markiert ein einmaliges Extremereignis --- oft eine Glitch-Entdeckung --- das zeitlich mit dem Chow-Bruch zusammenfallen, aber auch früher oder später auftreten kann.

Ein $\phi < 0{,}5$ zeigt, dass die Community nach einem Durchbruch die neue Technik verfeinert, statt grundlegend neue Möglichkeiten zu suchen. Ein $\phi > 1{,}0$ zeigt dagegen, dass der Durchbruch neue Glitch-Ketten oder Routing-Möglichkeiten freigelegt hat, die weitere schnelle Verbesserungen ermöglichen.

Für \textit{Portal~2}, \textit{Quake} und \textit{The Talos Principle} stimmt der Breakthrough-Zeitpunkt exakt mit dem Chow-Test-Strukturbruch aus F2 überein (jeweils WR~\#28, \#5 und \#23) --- ein Hinweis, dass der größte Einzelschritt der entscheidende Wendepunkt in der kumulierten Kurve ist. Bei \textit{Zelda: Ocarina of Time} liegt der maximale Einzelschritt (WR~\#4, Januar~2020) sechs~WRs vor dem Chow-Bruch (WR~\#10), was darauf hindeutet, dass ein früher Glitch-Fund das Routing fundamental veränderte, bevor sich die Verbesserungsgeschwindigkeit statistisch stabilisiert hatte.

\textit{Donkey Konga} ist der markanteste Ausreißer: $\phi = 23{,}97$ signalisiert eine Beschleunigung, nicht ein Plateau. Der Breakthrough (WR~\#9, Dezember~2024) liegt so kurz zurück, dass im Post-Segment kaum Konsolidierungszeit bestand --- die hohe Verbesserungsgeschwindigkeit danach spiegelt vor allem direkte Folgeverbesserungen wider. Alle übrigen auswertbaren Spiele ($n = 5$) zeigen ein klares Plateau ($\phi \leq 0{,}44$), wobei \textit{Zelda: OoT} ($\phi = 0{,}02$) und \textit{The Talos Principle} ($\phi = 0{,}02$) mit $\rho = -0{,}79$ ($p < 0{,}001$) die stärkste Abnahme der Post-Breakthrough-Verbesserungen aufweisen. Das Ergebnis zeigt: Durchbrüche setzen die Abflachungsrate nicht zurück, sondern leiten direkt eine Konsolidierungsphase ein.

\subsection*{F3b — Annäherung ans theoretische Minimum}

Der Asymptoten-Proxy schätzt das theoretische Minimum einer WR-Kurve über den Grenzwert des besten Sättigungsmodells. Der $R^2 \geq 0{,}70$-Schwellenwert schließt Spiele ohne stabile Sättigungskurve aus, damit keine unzuverlässigen Schätzungen einfließen. Dass 15 von 17 Spielen als \enquote{none\_detected} klassifiziert werden, ist kein Mangel der Methode, sondern ein inhaltliches Ergebnis: Diese Communities --- darunter jüngere Spiele wie \textit{Celeste} oder \textit{Hollow Knight} --- befinden sich noch in einer aktiven Entdeckungsphase, in der sich kein stabiler Grenzwert abzeichnet.

Für die zwei auswertbaren Spiele zeigt der Proxy, dass \textit{Half-Life~2} mit einem aktuellen WR von $2190{,}3$~s bereits 95,5~\% der möglichen Reduktion bis zum Modell-Floor ($1874{,}7$~s) erreicht hat. \textit{Pokémon Red/Blue} liegt bei 82,5~\%. Beide Spiele sind damit weit in die Sättigungsphase eingetreten.

Ein direkter Vergleich mit realen TAS-Zeiten war nicht möglich: Verlässliche TAS-Daten sind weder über eine öffentliche API verfügbar, noch lassen sie sich ohne Frame-genaue Videoanalyse zuverlässig erheben --- dabei ist zu berücksichtigen, dass TAS-Videos nicht immer vom tatsächlichen Spielstart beginnen, was die Zeitauswertung zusätzlich erschwert.

\subsection*{F3c — Lebensdauer}

Das Jahrzehntmuster --- lange Lebensdauer in den 2000ern, drastischer Rückgang ab 2010 --- passt zur historischen Entwicklung der Speedrunning-Community. Vor der Gründung von \textit{speedrun.com}~(2014) und der Popularisierung über Streaming-Plattformen gab es deutlich weniger aktive Speedrunner, sodass Weltrekorde länger unangefochten blieben. Der moderate Effekt ($\eta^2 \approx 0{,}044$) zeigt, dass das Genre zwar eine messbare Rolle spielt, aber bei weitem nicht der einzige Einflussfaktor ist. Community-Größe, Spielalter und Popularität auf Streaming-Plattformen sind wahrscheinlich stärkere Einflussfaktoren, lagen im Datensatz aber nicht vor.

Der hohe Gini-Koeffizient in allen Genres ($G = 0{,}65$–$0{,}92$) zeigt, dass Weltrekordverbesserungen sehr ungleich verteilt sind: Wenige Schlüsselentdeckungen (Glitches, neue Routen) machen den Großteil der Zeitersparnis aus, während die meisten späteren Rekorde nur winzige Sekundenbruchteile einsparen.

\subsection*{Einordnung in bestehende Forschung}

Erdil und Sevilla~\cite{erdil2023} zeigen anhand von 1.731~Weltrekorden aus 100~Spielen, dass WR-Verbesserungen Potenzgesetz-Mustern folgen und systematische Trendbrüche aufweisen. Diese Arbeit bestätigt den Strukturbruch-Befund (F2, Chow-Test für alle 17~Spiele signifikant) auf einem kleineren, aber detaillierter ausgewerteten Datensatz. Zusätzlich werden drei Fragen untersucht, die bei Erdil und Sevilla offen bleiben: Post-Breakthrough-Dynamik (F3a), Annäherung ans theoretische Minimum (F3b) und Genre-Vergleich der Lebensdauer (F3c). Ein direkter Ergebnisvergleich ist nicht möglich, da Spielsets, Zeiträume und Kategorien abweichen.

\section{Limitationen}

\textbf{Stichprobengröße:} Mit 17~Spielen ist der Datensatz klein. Verallgemeinerungen auf alle Genres sind daher nur eingeschränkt möglich.

\textbf{Kategorie any\%:} Die Beschränkung auf any\% schließt andere gängige Kategorien aus (z.\,B.~100\%, No-Major-Glitches). Die Dynamiken können sich je nach Kategorie erheblich unterscheiden.

\textbf{Plattformbias:} \textit{speedrun.com} hat einen Schwerpunkt auf PC-Spielen und englischsprachige Communities. Konsolen-Speedrunning und nicht-englischsprachige Spiele sind unterrepräsentiert.

\textbf{Fehlende Kontrollvariablen:} Community-Größe, Spieleranzahl und Medienaufmerksamkeit lagen nicht vor, obwohl sie Lebensdauer und Verbesserungsrate stark beeinflussen dürften.

\textbf{Kein TAS-Vergleich:} Ein direkter Vergleich mit TAS-Zeiten war nicht realisierbar. Keine öffentliche API liefert verlässliche TAS-Daten; die vorhandenen TAS-Videos erfordern Frame-genaue Auswertung des Startzeitpunkts, was im gegebenen Zeitrahmen nicht leistbar war. Der Asymptoten-Proxy ist daher nur ein Näherungswert für das theoretische Minimum.

\textbf{Kausalität:} Alle Ergebnisse zeigen Zusammenhänge, keine Ursachen. Aussagen wie „Glitches \emph{verursachen} Strukturbrüche" lassen sich aus diesen Daten nicht belegen.
